\section{Introduction}
\label{sec:intro}

Understanding how cells move from one state to another is a central problem in single-cell biology~\citep{farrell2018single,wagner2018single,schiebinger2019optimal,tanay2017scaling,griffiths2018using}. Processes such as development, differentiation, immune activation, tumor evolution, and cellular reprogramming are not simply collections of discrete cell types, but continuous population-level transitions through high-dimensional molecular state space~\citep{trapnell2014dynamics,haghverdi2016diffusion,setty2019characterization,cacchiarelli2018aligning}. The key questions are therefore dynamical, including which early states commit to particular fates, which intermediate states are transient but decisive, and which regulatory programs drive these transitions. Trajectory inference~\citep{lavenant2021towards,hashimoto2016learning} aims to answer these questions by reconstructing continuous cell-state evolution from time-course single-cell observations.

Transport- and flow-based trajectory methods provide a natural framework for modeling such dynamics~\citep{schiebinger2019optimal, peyre2019computational,lavenant2021towards}. Since time-course single-cell data provide population-level snapshots rather than paired observations of the same cells~\citep{schiebinger2019optimal, weinreb2018fundamental}, these methods must infer how marginal distributions are coupled across time. Most existing formulations operate in the Euclidean representation space, where couplings are estimated from Euclidean endpoint costs and continuous paths or velocity fields are regularized by Euclidean kinetic energy, action, or smoothness~\citep{neklyudov2023action,tong2023conditional,tong2020trajectorynet,huguet2022manifold}. This Euclidean geometry is convenient, but it is not necessarily aligned with biological progress. Two cells can be close in expression space while lying at different developmental stages or fate branches, whereas biologically plausible motion may follow curved and locally anisotropic directions along the cell-state manifold~\citep{moon2019visualizing,wolf2019paga}.

This mismatch makes trajectory inference from destructive snapshots fundamentally ambiguous~\citep{weinreb2018fundamental,tritschler2019concepts}: the cross-time coupling is not observed and is generally not identifiable from marginal distributions alone. A coupling may satisfy the marginal constraints while connecting cells across incompatible developmental programs, stalled states, or fate branches. For example, in mouse reprogramming~\citep{schiebinger2019optimal}, differences in growth rates between cell types can make marginal matching misleading, coupling apoptotic stromal cells to rapidly expanding iPSCs. % PAGECUT: sentence tightened for the ICLR page limit. Original:
% "Human iPSC reprogramming~\citep{liu2020reprogramming} shows the same issue at finer granularity, that cells collected on the same day can contain heterogeneous primed-like, naive-like, and trophectoderm-like intermediates, so clock-time adjacency alone does not determine which cells should be coupled."
Human iPSC reprogramming~\citep{liu2020reprogramming} shows the same issue at finer granularity: cells collected on the same day span primed-like, naive-like, and trophectoderm-like intermediates, so clock-time adjacency alone does not determine which cells should be coupled.
Auxiliary measurements such as RNA velocity, lineage tracing, or metabolic labeling~\citep{la2018rna,bergen2020generalizing,lange2022cellrank} can help, but require additional experiments and are unavailable in many datasets. Existing Euclidean or support-based regularizers can encourage short, smooth, or data-supported paths~\citep{saelens2019comparison}, but they do not directly encode which local directions of motion are developmentally admissible.

We propose PACE, a geometry-aware trajectory-inference framework for selecting snapshot-consistent continuous transport dynamics from unpaired time-course snapshots.
The intuition is that, although true cell identities and intermediate paths are unobserved, local spatiotemporal neighborhoods still provide a weak but useful prior on admissible directions, since motion along locally supported tangent directions is more plausible than motion orthogonal to the observed manifold structure.
PACE therefore replaces Euclidean transport costs with a Riemannian path action induced by a state- and time-dependent anisotropic metric, which assigns lower cost to locally supported tangent motion and higher cost to normal motion.
Specifically, PACE estimates local tangent \emph{directions} from spatiotemporal neighborhoods of observed cells and uses this metric to define an anisotropic bridge transport problem, where the cost of coupling two cells is the minimum Riemannian action of an endpoint-conditioned path.
PACE then alternates between refining OT couplings under this path-action cost and fitting endpoint-preserving neural bridges, allowing local geometry to influence both which cells are coupled and how they move between snapshots.
Finally, the learned bridge dynamics are distilled into a continuous-time population velocity field over cellular states.

Our contributions are threefold:
% PAGECUT: blank lines between items in this list were removed for the ICLR
% page limit (spacing-only change; the NeurIPS source had blank lines between items).
\begin{enumerate}
    \item We introduce a time- and state-dependent Riemannian metric based on local spatiotemporal tangent-space projections, providing a geometry-aware path-action cost for selecting couplings between adjacent destructive snapshots beyond Euclidean endpoint proximity.
    \item We develop an iterative bridge-transport procedure that alternates between Riemannian coupling refinement and endpoint-preserving neural bridge fitting, allowing cross-time correspondences and interpolant paths to be selected jointly under the same geometry-aware action.
\item 
We distill endpoint-conditioned bridges into a global continuous-time population velocity field and evaluate PACE on seven datasets covering nine held-out reconstruction experiments. 
PACE achieves the strongest overall reconstruction performance, reducing MMD, $\mathcal{W}_1$, and $\mathcal{W}_2$ by 23.7\% on average relative to the strongest competing baseline, improves RNA-velocity alignment on an embryoid body differentiation benchmark, and shows consistent gains across component ablations.
\end{enumerate}
